% ============================================================
% content/answers-short.tex —— 2023 湖南卷物理 · 简明答案速查
% ============================================================

\begin{gknotice}
  选择题答案：1--A，2--B，3--C，4--B，5--D，6--D，7--BC，8--AD，9--AC，10--CD
\end{gknotice}

\section*{一、选择题}

\noindent
\begin{tabular}{@{}llllllllll@{}}
\toprule
题号 & 1 & 2 & 3 & 4 & 5 & 6 & 7 & 8 & 9\\
\midrule
答案 & A & B & C & B & D & D & BC & AD & AC\\
\bottomrule
\end{tabular}

\smallskip

\noindent
第 10 题：CD

\section*{二、非选择题}

\noindent
\textbf{11.} (3) $T=\dfrac{t_{0}}{10}$；\quad(5) 线性的；\quad(6) A；\quad
(7) 弹簧非轻质（其质量不可忽略），或手机内部的磁体、导体使振动受到电磁阻尼（任选一条，合理即可）

\smallskip

\noindent
\textbf{12.} (1) $1000$；\quad(2) $R_{F}=\dfrac{R_{1}R_{3}}{R_{2}}$；\quad
(3) 图线为过原点的直线；\quad(4) $1.75\times10^{-2}$；\quad(5) $>$

\smallskip

\noindent
\textbf{13.} (1) $p_{1}=\dfrac{V_{0}}{V_{0}+V_{1}}p_{0}$；\quad
(2) $\Delta F=\left[1-\left(\dfrac{V_{0}}{V_{0}+V_{1}}\right)^{n}\right]p_{0}S$

\smallskip

\noindent
\textbf{14.} (1) $v_{0}=\dfrac{2mgR\sin\theta}{B^{2}L^{2}}$；\quad(2) $a_{0}=2g\sin\theta$；\quad
(3) $v=gt_{0}\sin\theta+\dfrac{mgR\sin\theta}{B^{2}L^{2}}$，\quad
$\Delta x=\dfrac{2m^{2}R^{2}g\sin\theta}{B^{4}L^{4}}$

\smallskip

\noindent
\textbf{15.} (1) $v_{1}=\sqrt{\dfrac{2m^{2}gb}{m_{0}(m_{0}+m)}}$，凹槽向右移动 $x_{1}=\dfrac{ma}{m_{0}+m}$；\quad
(2) $\dfrac{\left(x-\frac{m}{m_{0}+m}a\right)^{2}}{\left(\frac{am_{0}}{m_{0}+m}\right)^{2}}+\dfrac{y^{2}}{b^{2}}=1$（$y\le0$）；\quad
(3) $v=\sqrt{\dfrac{4gb^{2}}{a+3b}}$
